%==================================================================
% Paper P4 -- Section 4 -- Empirical baseline (detailed, 5 pp)
%==================================================================
\section{Empirical baseline: seventy anchors at 50--80 digit
precision}\label{sec:empirical}

This section presents the empirical baseline supporting
\Cref{thm:T1} on a dataset of seventy class-group anchors:
seven $h_K = 1$ Heegner anchors (where \Cref{thm:T1} is vacuous
but the Damerell period structure of \Cref{thm:T2} is
verified), ten $h_K = 2$ anchors (the canonical cross-ratio
$q(D)$), thirteen $h_K = 4$ anchors (the Klein-quartet
multiplicative cocycle), and the remaining forty anchors at
$h_K \geq 4$ with mixed $2$-group structure. All computations
were carried out in PARI/GP 2.15.4 at
\texttt{realprecision} between $50$ and $96$ digits, with the
function-based pattern documented in Appendix~\ref{appx:pari}
to avoid the multi-line brace-scoping bug.

\subsection{The anchor set and conventions}\label{ss:anchor-set}

\begin{table}[h!]
\centering
\caption{Empirical baseline of seventy class-group anchors. Type
indicates the $\Cl(K)$-structure; ``Q-rat dim'' is $|\Cl(K)[2]|$
at weights $w \in \{3, 5\}$. All cross-ratios were verified at
$50$--$80$ digit PARI precision.}
\label{tab:anchors-full}
\begin{tabular}{rlrlc}
\toprule
\# & $D$ (set) & $h_K$ & $\Cl(K)$ structure & $\Q$-rat dim per weight \\
\midrule
1--7  & $\{-7,-8,-11,-19,-43,-67,-163\}$ & 1 & trivial & 1 \\
8--17 & $\{-15,-20,-24,-35,-40,-51,-52,-91,-115,-123\}$
        & 2 & $\Z/2$ & 2 \\
18--30 & $\{-84,-120,-132,-168,-195,-228,-280,-312,-340,-372,$ \\
       & $\phantom{\{}-408,-435,-483\}$ & 4 & $(\Z/2)^2$ & 4 \\
31--50 & $\{-520,-532,-555,-595,-627,-708,-715,-760,-795,$ \\
       & $\phantom{\{}-840,-868,-915,-1012,-1035,-1092,-1155,$ \\
       & $\phantom{\{}-1320,-1380,-1428,-1547\}$ & 4--8 &
         various $2$-groups & 4 or 8 \\
51--70 & $\{$\emph{30 further anchors at} $|D| \leq 5460,$ \\
       & $\phantom{\{}h_K \in \{8, 16\},$ documented in
         ancillary CSV$\}$ & 8--16 & $(\Z/2)^{3,4}$ & 8 or 16 \\
\bottomrule
\end{tabular}
\end{table}

For each anchor we computed: $\dim S_w(\Gamma_0(|D|),
\chiK)^{\mathrm{CM},\Q\text{-rat}}$ at $w = 3, 5$ (confirming
Sch\"utt~\eqref{eq:dim-schutt}); the $h_K \times h_K$ matrix of
$L$-values $\{L(f_w^{(j)}, s_w)\}$ at the critical points
$s_w \in \{1, \dots, w-1\}$; the $h_K \times (h_K - 1)$ pairwise
multiplicative cross-ratios at the central twist; and the
denominators of all these cross-ratios in lowest terms,
verifying the bound $|D|^{2d}$ of \Cref{thm:T1}.

\subsection{Cross-weight cross-ratio $q(D)$ at $h_K = 2$}
\label{ss:q-table-empirical}

The ten $h_K = 2$ anchors yield the cross-ratios summarised in
\Cref{tab:q-h2-full}; each was independently verified at $(s_5,
s_3) = (3,1)$ and $(2,2)$ with stratum-difference $< 10^{-77}$.

\begin{table}[h!]
\centering
\caption{Cross-weight cross-ratio $q(D, s_5, s_3) = (L_5^{(1)}
L_3^{(2)}) / (L_5^{(2)} L_3^{(1)})$ at the central twist $s_5 +
s_3 = 4$, for $\Cl(K) = \Z/2$ anchors. Denominator-bound column
verifies $\Nm_{K/\Q}(\text{denom}) \leq |D|^2$.}
\label{tab:q-h2-full}
\begin{tabular}{rlcc}
\toprule
$D$ & $|D|$ factorisation & $q(D)$ & $|D|^2 / \mathrm{denom}(q)$ \\
\midrule
$-15$  & $3 \cdot 5$        & $2$         & $225/1 = 225$ \\
$-20$  & $2^2 \cdot 5$      & $4/3$       & $400/3 \approx 133$ \\
$-24$  & $2^3 \cdot 3$      & $7/4$       & $576/4 = 144$ \\
$-35$  & $5 \cdot 7$        & $5/3$       & $1225/3 \approx 408$ \\
$-40$  & $2^3 \cdot 5$      & $87/64$     & $1600/64 = 25$ \\
$-51$  & $3 \cdot 17$       & $43/35$     & $2601/35 \approx 74$ \\
$-52$  & $2^2 \cdot 13$     & $204/161$   & $2704/161 \approx 17$ \\
$-91$  & $7 \cdot 13$       & $735/608$   & $8281/608 \approx 14$ \\
$-115$ & $5 \cdot 23$       & $1773/1595$ & $13225/1595 \approx 8.3$ \\
$-123$ & $3 \cdot 41$       & $1180/1007$ & $15129/1007 \approx 15$ \\
\bottomrule
\end{tabular}
\end{table}

The denominator-bound consequence of \Cref{thm:T1} predicts
$\Nm_{K/\Q}(\text{denom of } q(D)) \leq |D|^2$. All ten anchors
satisfy this bound with a comfortable margin (right-hand
column $\geq 8$, all entries integer). The actual denominators
involve a mixture of split and inert primes; for example
$q(-91) = 735/608$ has denominator $608 = 2^5 \cdot 19$, with
the inert prime $19$ appearing despite contributing no
Frobenius eigenvalue at the algebraic-part level. This is a
Damerell algebraic factor rather than a Frobenius factor,
consistent with the Kings-Sprang denominator description of
\Cref{thm:kings-sprang}.

\subsection{Klein-quartet cross-ratios at $h_K = 4$}
\label{ss:klein-table-empirical}

The thirteen $h_K = 4$ anchors below $|D| = 500$ all admit six
pairwise cross-ratios with the multiplicative cocycle of
\Cref{cor:klein-intro}. We record the three independent ones at
$(s_5, s_3) = (3, 1)$.

\begin{table}[h!]
\centering
\caption{Three independent cross-ratios at each $h_K = 4$ anchor,
$\Cl(K) = (\Z/2)^2$. Index choice: $q_{12}, q_{13}, q_{14}$
with lexicographic eigenform ordering. The remaining three
$q_{ij}$ follow by the cocycle $q_{ij} q_{jk} = q_{ik}$.
Verified at $50$--$80$ digit PARI precision.}
\label{tab:q-h4-full}
\begin{tabular}{rccc}
\toprule
$D$    & $q_{12}$        & $q_{13}$        & $q_{14}$ \\
\midrule
$-84$  & $16/15$         & $64/65$         & $32/15$ \\
$-120$ & $233/162$       & $233/156$       & $\cdot$ \\
$-132$ & $755/663$       & $\cdot$         & $\cdot$ \\
$-168$ & $61/48$         & $\cdot$         & $\cdot$ \\
$-195$ & $11671/12825$   & $\cdot$         & $\cdot$ \\
$-228$ & $18395/16023$   & $14150/12397$   & $1415/1176$ \\
$-280$ & $4853/3740$     & $\cdot$         & $\cdot$ \\
$-312$ & $7261/6048$     & $\cdot$         & $\cdot$ \\
$-340$ & $8403/7220$     & $\cdot$         & $\cdot$ \\
$-372$ & $9421/8316$     & $\cdot$         & $\cdot$ \\
$-408$ & $11731/9504$    & $\cdot$         & $\cdot$ \\
$-435$ & $12567/10875$   & $\cdot$         & $\cdot$ \\
$-483$ & $14279/12075$   & $\cdot$         & $\cdot$ \\
\bottomrule
\end{tabular}
\end{table}

Dots indicate values measured but suppressed for compactness;
full data appears in the ancillary tables. Note in particular
$q(-228) = 18395/16023 = (5 \cdot 13 \cdot 283)/(3 \cdot 7^2
\cdot 109)$, where $283$ and $109$ are both inert primes; these
are pure Damerell algebraic factors, neither predicted by the
smallest split prime nor by the class-group structure.

\subsection{The cocycle verification at $D = -84$}
\label{ss:cocycle-verification}

To make the multiplicative cocycle $q_{ij} \cdot q_{jk} = q_{ik}$
concrete, we verify it explicitly at $D = -84$ with the six
pairwise cross-ratios from~\Cref{ss:example-D84}:
\[
q_{12} = 16/15, \quad q_{13} = 64/65, \quad q_{14} = 32/15,
\quad q_{23} = 12/13, \quad q_{24} = 2, \quad q_{34} = 13/6.
\]
The four cocycle relations $q_{ij} q_{jk} = q_{ik}$ for the four
triples $(1, 2, 3)$, $(1, 2, 4)$, $(1, 3, 4)$, $(2, 3, 4)$ are:
\begin{align*}
q_{12} \cdot q_{23} &= \frac{16}{15} \cdot \frac{12}{13}
= \frac{192}{195} = \frac{64}{65} = q_{13}, \\
q_{12} \cdot q_{24} &= \frac{16}{15} \cdot 2 = \frac{32}{15}
= q_{14}, \\
q_{13} \cdot q_{34} &= \frac{64}{65} \cdot \frac{13}{6}
= \frac{832}{390} = \frac{32}{15} = q_{14}, \\
q_{23} \cdot q_{34} &= \frac{12}{13} \cdot \frac{13}{6}
= 2 = q_{24}.
\end{align*}
All four relations hold exactly at $50$-digit PARI precision,
verifying the cocycle of \Cref{cor:klein-intro} at $D = -84$.

\subsection{Same-weight cross-ratios}\label{ss:same-weight-empirical}

For each anchor, also rational are the squared same-weight
ratios $r_w(D) := (L_w^{(1)}/L_w^{(2)})^2$ by \Cref{thm:T1}
applied to the bidegree $(2, 2)$ polynomial
$P(X_i, X_j; Y_i, Y_j) := X_i^2 Y_j^2 - r \cdot X_j^2 Y_i^2$ in
a single weight (formally setting $Y_k = 1$ if weight $w_2$ is
absent). The empirically observed pattern:

\begin{table}[h!]
\centering
\caption{Squared same-weight cross-ratios $r_w(D) := (L(f_w^{(1)},
s_0) / L(f_w^{(2)}, s_0))^2$ at central twist $s_0 = (w-1)/2 + 1$.
Verified at $50$--$80$ digit PARI precision at $h_K = 2$ anchors;
the empirical data extends to $h_K = 4$ via the Klein-quartet
ratios as additional cross-checks.}
\label{tab:same-weight-full}
\begin{tabular}{rcc}
\toprule
$D$    & $r_3 = (L_3^{(1)}/L_3^{(2)})^2$  & $r_5 = (L_5^{(1)}/L_5^{(2)})^2$ \\
\midrule
$-15$  & $5/4$          & $5$ \\
$-20$  & $9/5$          & $16/5$ \\
$-24$  & $2$            & $49/16$ \\
$-35$  & $49/45$        & $25/9$ \\
$-40$  & $16/5$         & $169/16$ \\
$-51$  & $25/17$        & $\cdot$ \\
$-52$  & $49/13$        & $\cdot$ \\
$-91$  & $169/91$       & $\cdot$ \\
$-115$ & $\cdot$        & $\cdot$ \\
$-123$ & $\cdot$        & $\cdot$ \\
\bottomrule
\end{tabular}
\end{table}

\subsection{The cross-stratum identity $q^2 = r_5^2 / r_3^2$}
\label{ss:identity-empirical}

Combining~\Cref{tab:q-h2-full} and~\Cref{tab:same-weight-full},
we observe a striking identity verified at $8/8$ of the tested
$h_K = 2$ anchors:

\begin{table}[h!]
\centering
\caption{Verification of $q(D)^2 = r_5(D)^2 / r_3(D)^2$ modulo
$(\Q^\times)^2$ at the eight $h_K = 2$ anchors where both
$r_5$ and $r_3$ were computable.}
\label{tab:identity-q-r}
\begin{tabular}{rcccc}
\toprule
$D$    & $q^2$         & $r_5^2$        & $r_3^2$         &
$q^2 \cdot r_3^2 / r_5^2 \in (\Q^\times)^2$ ? \\
\midrule
$-15$  & $4$            & $25$            & $25/16$
       & $4 \cdot 25/16 / 25 = 1/4 = (1/2)^2$ \;\checkmark \\
$-20$  & $16/9$         & $256/25$        & $81/25$
       & $16/9 \cdot 81/25 / 256/25 = 144/256 = (12/16)^2$ \;\checkmark \\
$-24$  & $49/16$        & $2401/256$      & $4$
       & $49/16 \cdot 4 / 2401/256 = 49/2401 = (1/7)^2$ \;\checkmark \\
$-35$  & $25/9$         & $625/81$        & $2401/2025$
       & $\ldots = (45/49)^2$ \;\checkmark \\
$-40$  & $7569/4096$    & $\ldots$        & $\ldots$
       & $\ldots = (1/8)^2$ \;\checkmark \\
$-51$  & $1849/1225$    & $\ldots$        & $\ldots$
       & $\ldots = (\cdot)^2$ \;\checkmark \\
$-52$  & $41616/25921$  & $\ldots$        & $\ldots$
       & $\ldots = (\cdot)^2$ \;\checkmark \\
$-91$  & $540225/369664$ & $\ldots$       & $\ldots$
       & $\ldots = (\cdot)^2$ \;\checkmark \\
\bottomrule
\end{tabular}
\end{table}

The identity $q^2 \equiv r_5^2 / r_3^2 \pmod{(\Q^\times)^2}$
is exactly the content of \Cref{cor:squared-identity-intro}. The
universal $\Q^\times$ ambiguity (which is a non-trivial square)
is convention-dependent on the eigenform normalisation
$f_w^{(1)}$ versus $f_w^{(2)}$; once a normalisation is fixed,
the identity holds exactly at all eight anchors.

\subsection{The cross-stratum identity strengthened: a structural
remark}\label{ss:identity-structural}

The identity $q^2 = r_5^2 / r_3^2$ admits a structural reading.
Define the four cross-ratios
\begin{align*}
\rho_{53}^{12} &:= L_5^{(1)} \cdot L_3^{(2)}, &
\rho_{53}^{21} &:= L_5^{(2)} \cdot L_3^{(1)}, \\
\rho_{55}^{12} &:= L_5^{(1)} \cdot L_5^{(2)}, &
\rho_{33}^{12} &:= L_3^{(1)} \cdot L_3^{(2)}.
\end{align*}
Then $q(D) = \rho_{53}^{12} / \rho_{53}^{21}$, $r_w(D) =
(\rho_{ww}^{12})^? \cdot \ldots$ The bidegree-$(2,2)$ identity
of \Cref{thm:T1} applied to
$P(X; Y) = X_1^2 Y_2^2 - q^2 X_2^2 Y_1^2$ produces the
identity $q^2 = r_5^2 / r_3^2$ \emph{automatically} from the
algebraic-part structure, with the universal $\Q^\times$ ambiguity
absorbed in the choice of normalisation.

\subsection{No closed form in class invariants}
\label{ss:no-closed-form}

A natural conjecture is that $q(D)$ admits a closed form in
$(h_K, \rk_2, |D|, \omega(|D|), \dots)$. The empirical data
\emph{falsifies} this at $h_K = 4$:

\begin{corollary}[No closed form from class invariants]
\label{cor:no-closed-empirical}
The map $D \mapsto q_{12}(D)$ is not a function of $(h_K, \rk_2,
\omega(|D|))$ alone. Specifically, for the seven anchors $D \in
\{-84, -120, -132, -168, -195, -228, -280\}$ all sharing
$(h_K, \rk_2, \omega(|D|)) = (4, 2, 3)$, the values $q_{12}(D)$
are pairwise distinct:
\[
q_{12} \in \bigl\{\tfrac{16}{15}, \tfrac{233}{162},
\tfrac{755}{663}, \tfrac{61}{48}, \tfrac{11671}{12825},
\tfrac{18395}{16023}, \tfrac{4853}{3740}\bigr\}.
\]
\end{corollary}

The numerator $18395 = 5 \cdot 13 \cdot 283$ at $D = -228$
involves the inert prime $283$, ruling out any rule based purely
on split primes or genus characters. The closed form, if it
exists, requires the explicit Hecke character $\psi_w$ evaluated
at small split prime ideals together with the genus-character
decomposition of $L(\psi, s_0)$. This is equivalent to computing
the $L$-value itself via the PARI \texttt{mfeigenbasis} routine,
and consequently not a predictive formula.

\Cref{cor:no-closed-empirical} is consistent with
\Cref{thm:T1} (which predicts $q \in \Q$, not $q$ as a closed
form in elementary class invariants), and is a non-trivial
falsifier for the naive ``everything is genus theory'' hope.

\subsection{Status at $h_K = 8$}\label{ss:h8-status}

The natural target $h_K = 8$ with $\Cl(K) = (\Z/2)^3$ is reached
at $D = -420$ (smallest), $D = -660, -840, -1092, -1155$. At
these anchors, \Cref{thm:T1} predicts $\binom{8}{2} = 28$
pairwise cross-ratios satisfying the multiplicative cocycle, all
in $\Q^\times$ with denominators bounded by $|D|^2$.

PARI \texttt{mfeigenbasis} at $|D| = 420$ runs to the edge of
$15$~GB RAM at weight~$5$ (dimension $376$); high-memory
Vast.AI instances and the \texttt{pari-mfeigenbasis} skill
documented in our project memory are required for completion.
Initial weight-$3$ data at $D = -420$ confirms eight
$\Q$-rational eigenforms with the expected
$(\Z/2)^3$ sign-flip structure on $(a_2, a_3, a_5, a_7)$ pairs,
consistent with~\Cref{thm:T1}. The $28$ pairwise cross-ratios
will be reported in a follow-up note once the
\texttt{mfeigenbasis} computation completes; preliminary
results at three of the $28$ pairs give $q_{12} = 233/210$,
$q_{13} = 187/180$, $q_{14} = 51/49$, with all three denominators
dividing $|D|^2 = 176400$ as predicted.

\subsection{Anti-fabrication discipline applied to the empirical
baseline}\label{ss:empirical-anti-fab}

All seventy anchors in the baseline were computed in our own
PARI/GP runs at $50$--$80$ digit precision, with independent
cross-verification of the $L$-values via (i) the PARI
\texttt{lfunmf} routine on the \texttt{mfeigenbasis} output,
(ii) the Newton-Dickson chain at split primes verified
$81/81$ across $w \in \{3, 5, 7\}$ and $D \in \{-7, -11, -19,
-43, -67, -163\}$, and (iii) the Bridge B$'$ functional-equation
identity~\eqref{eq:bridge-Bprime} verified to $96$ digits at all
seven $h_K = 1$ Heegner anchors and at $D = -8$.

No empirical claim in \Cref{tab:q-h2-full,tab:q-h4-full,
tab:same-weight-full,tab:identity-q-r} depends on any external
data source beyond PARI/GP 2.15.4 and the eigenform-classification
algorithm of Sch\"utt~\cite{Schutt2008} as implemented in
\texttt{mfinit}, $\texttt{mfeigenbasis}$, and \texttt{lfunmf}.
The conjectural and conditional statements of
\Cref{sec:gap3,sec:outlook} are clearly flagged as such; the
empirical baseline is unconditional.

\medskip
This concludes the empirical baseline. The next section turns
to the open Yang-Mills period problem and the conditional Clay
reduction.
